%%
%% fall-lecture18.tex
%% 
%% Made by Alex Nelson <pqnelson@gmail.com>
%% Login   <alex@lisp>
%% 
%% Started on  2026-11-08T13:01:20-0800
%% Last update 2026-11-08T13:01:20-0800
%% 

\lecture{}

Properties of the Euler class.

\begin{node}[Property 1]
Naturality: $f^{*}e(\xi)=e(f^{*}\xi)$ [this follows by naturality of $U$].
\end{node}

\begin{corollary}
If $\xi$ is a trivial $\RR^{n}$-bundle (when $n\geq1$), then
$e(\xi)=0$ by similar reasoning as the Stiefel--Whitney class being trivial.
\end{corollary}

\begin{node}[Property 2]
Reversing the orientation of $\xi$ sends $U\mapsto-U$, and so
$e(\xi)\mapsto-e(\xi)$. 
\end{node}

\begin{remark}
As an immediate consequence to this, the uniqueness of integral Thom
class is not guaranteed since we can always choose $-U$ or $+U$.
\end{remark}

\begin{node}[Property 3]
If $\xi$ has odd rank, then $e(\xi)+e(\xi)=0$.
\end{node}

\begin{proof}[First proof]
Let $\xi\xrightarrow{\iso}\xi$ sending $(x,v)\mapsto(x,-v)$ which
flips the orientation on fibers. This gives us $e(\xi)=-e(\xi)$.
\end{proof}

\begin{proof}[Second proof]
Use integral version of Thom isomorphism, this will give us
$U\smile U=(-1)^{n^{2}}U\smile U=-U\smile U$ since $n$ is odd.
The result follows immediately from this.
\end{proof}

\begin{node}[Property 4]
We see $w_{n}(\xi)=e(\xi)\pmod{2}$.
\end{node}

\begin{node}[Property 5]
Assume $\xi$, $\xi'$ are rank-$m$ and rank-$n$ bundles over the same
base $B$. Then
\begin{equation}
e(\xi\oplus\xi')=e(\xi)\smile e(\xi'),
\end{equation}
and
\begin{equation}
e(\xi\times\xi')=e(\xi)\times e(\xi').
\end{equation}
The proof is basically the same as in Property~4, from
$e(\xi\times\xi')=e(\xi)\times e(\xi')$ we pullback
$\Delta^{*}(\xi\times\xi')$ to prove the first part of the claim.
\end{node}

\begin{node}[Property 6]
If $\xi$ has a nowhere zero section, then $e(\xi)=0$.
\end{node}

\begin{proof}
Assume $\xi$ has a Euclidean metric (as an extra assumption). [This
  covers the paracompact $B$ case.]
Then there is an $\eta$ such that $\xi=\eta\oplus\eta^{\perp}$. Then
\begin{subequations}
  \begin{align}
e(\xi) &= e(\eta)\smile e(\eta^{\perp})\\
&= 0\smile e(\eta^{\perp})\\
&= 0.
  \end{align}
\end{subequations}
Hence the claim.
\end{proof}

\begin{theorem}[Very important!!!]
If $M$ is a smooth manifold and $\xi=\T M$ is its tangent bundle,
then $\langle e(\xi),[M]\rangle=\chi(M)$ is the Euler characteristic
of $M$.
\end{theorem}

The left-hand side $\langle e(\xi),[M]\rangle$ describes something
about the geometry of the bundle, the right-hand side $\chi(M)$
describes something about the topology of $B$ (its base space).